\documentclass[11pt]{article}
\usepackage{explorer}
\newcommand{\ExplorerLanguage}{FR}
\newcommand{\ExplorerNumber}{11}
\newcommand{\ExplorerKind}{Exercices progressifs}
\begin{document}
\ExplorerTitle{Lycée · optimisation et olympiades}{Inégalités et cas d’égalité}
\ExplorerSection{Exercice 1}
Pour \(x>0\), prouver \(x+\frac{1}{x}\geq 2\) et déterminer l’égalité.\par
\vspace{5pt}
\ExplorerSection{Exercice 2}
Déterminer le minimum de \(x^{2}-6x+13\) sur \(\mathbb{R}\).\par
\vspace{5pt}
\ExplorerSection{Exercice 3}
Un rectangle a pour périmètre 24 m. Quelle est son aire maximale ?\par
\vspace{5pt}
\ExplorerSection{Exercice 4}
Pour a,b réels, prouver \(a^{2}+b^{2}\geq 2ab\). Faut-il supposer a et b positifs ?\par
\vspace{5pt}
\end{document}
